\section*{Evaluation settings and remaining reproducibility records}

\subsection*{Validity checks and aligned distances}

The configured PoseBusters assessment checks loading, sanitization, InChI conversion,
connectivity, radicals, molecular formula, bond connectivity, tetrahedral chirality,
and double-bond stereochemistry. Distance-geometry checks ignore hydrogens and use
configured tolerances of 0.25 for bond lengths and angles and 0.3 for internal clashes.
The bounds matrix enables 1--5 bounds, van der Waals scaling, and triangle smoothing,
and disables the macrocycle 1--4 option. Aromatic five- and six-membered rings and
trigonal carbon--carbon double bonds have a maximum deviation from planarity of
0.25 \AA{}. The energy-ratio threshold is 100, with 50 internally generated comparison
conformers. This validity calculation is separate from the common-hydrogen,
unoptimized MMFF94s energy analysis reported in the Results. The inspected configuration
is archived as \texttt{posebusters-config.json}; no receptor-contact tests are included.

RMSD calculations use RDKit's \texttt{GetBestRMS} after removing hydrogens, without
explicitly overriding its symmetry-search defaults. CASF comparisons that raise an
alignment exception attempt an index-based atom mapping; failures of both procedures
are unmeasurable. The PLINDER matching routine instead records failed alignments as
undefined. All nearest-reference and nearest-generated distances used in the reported
PLINDER metrics are finite, which does not establish that every pairwise comparison
succeeded. Greedy cluster representatives remain fixed in saved conformer order.

\textbf{[TODO: archive the executed RDKit and PoseBusters versions and configurations
for each historical evaluation, including the energy-check random seed and any
runtime overrides. Count index-mapping fallbacks and failed pairwise alignments,
and report the symmetry-search limits and rotatable-bond descriptor option for
those versions. Do not infer these from the currently installed environment.]}

\subsection*{Molecule-level displays}

The trajectories in main Figure 6 identify the union of
molecules with the largest MAT-R or MAT-P for any displayed method; the other
molecules are grey. For Supporting Figure S6, each molecule's contribution to the
validity change is its change in pass percentage divided by the sum of these changes
across molecules. This decomposes the net change in the unweighted mean pass rate,
not the pooled rate in the main comparison. Negative contributions remain negative.
